\documentclass[10pt,a4paper]{amsart}
\usepackage[margin=19mm]{geometry}
\usepackage{amsmath,amssymb,mathtools}
\usepackage[scr=rsfs,cal=euler]{mathalfa}
\usepackage{xfrac}
\usepackage[unicode,hidelinks,bookmarks=false]{hyperref}
\newcommand{\ie}{i.e.\ }
\newcommand{\cf}{cf.\ }
\newcommand{\resp}{respectively\ }
\newcommand{\Subsection}{\subsection}
\providecommand{\nobreakdash}{\nobreak\hbox{-}}
\setlength{\emergencystretch}{2em}
\multlinegap0pt
\usepackage{stmaryrd}
\theoremstyle{plain}
\newtheorem{thm}{Theorem}[section]
\newtheorem{theorem}[thm]{Theorem}
\newtheorem{cor}[thm]{Corollary}
\newtheorem{corollary}[thm]{Corollary}
\newtheorem{lem}[thm]{Lemma}
\newtheorem{lemma}[thm]{Lemma}
\newtheorem{claim}[thm]{Claim}
\newtheorem{prop}[thm]{Proposition}
\newtheorem{proposition}[thm]{Proposition}
\newtheorem{defn}[thm]{Definition}
\newtheorem{definition}[thm]{Definition}
\newtheorem{exmp}[thm]{Example}
\newtheorem{example}[thm]{Example}
\newtheorem{rem}[thm]{Remark}
\newtheorem{remark}[thm]{Remark}
\newcommand{\getl}{\mathsf{getL}}
\newcommand{\putr}{\mathsf{putR}}
\newcommand{\getr}{\mathsf{getR}}
\newcommand{\putl}{\mathsf{putL}}
\newcommand{\corner}[1]{{\text{}^\ulcorner_\llcorner\!{#1}\!_\lrcorner^\urcorner}}
\newcommand{\asn}{\downarrow}
\newcommand{\seq}{,\!\ldots\!,}
\title{A Simple Categorical Calculus of Interacting Processes}
\author{Chad Nester and Niels Voorneveld}
\date{Source: arXiv:2603.18321v2 (CC BY 4.0); the authors' own native \LaTeX{} source, set here in the editorial AMS \texttt{amsart} wrapper (the source's own class is the Entics \texttt{entics} class, which is not redistributed).}
\begin{document}
\maketitle
\noindent\textit{Editorial scope.} The counted claim of this unit is the article's Lemma 3.7 (source label \texttt{lem:substitution-coherent}), the substitution coherence of the interpretation of the term calculus, delivered together with the authors' complete original proof. The statement and the proof are the only retained passages, and every one of their characters is the authors' own; the proof is delivered as the single primary proof body exactly as the authors' proof environment encloses it, including the two lengthy aligned computations that make up the value case and the nested let case. Printed numbering: the authors' theorem-like environments are declared inside the Entics class that the article uses, which is not present in the runtime and is not redistributed, so the wrapper restates the same declarations, with the lemma sharing the counter of the theorem environment and that counter being numbered by section exactly as the class does, and advances the section and theorem counters so that the delivered PDF prints the counted claim as Lemma~3.7, exactly as the published article does. Omitted: the rest of the article, that is its title block, abstract, introduction, the sections on the free cornering and the term calculus that precede the lemma, and the sections that follow it; nothing of the counted statement or of the counted proof is omitted, and the retained unit is the maximal source-local passage of the article that is closed under its own cross-references. Class substitution: the source class is the publisher's Entics \texttt{entics} class together with the author's own \texttt{enticsmacro} style file; neither file is present in this package and neither is redistributed. Reference handling: no retained passage contains a citation; the article's bibliography is a separate file that is kept verbatim in the eprint directory and is not redistributed, and no reference entry is transcribed, delivered or invented. No mathematics is written, repaired or re-derived here.

\setcounter{section}{3}
\setcounter{thm}{6}
\begin{lemma}\label{lem:substitution-coherent}
  $\llbracket \textsf{let } x_1 \seq x_n \asn ([v_1] \mid \cdots \mid [v_n]) \textsf{ in } t \rrbracket = \llbracket t[v_1 \seq v_n / x_1 \seq x_n]\rrbracket$ whenever this makes sense.
\end{lemma}

\begin{proof}
  We proceed by structural induction on $t$. The cases are as follows:
  \begin{itemize}
  \item If $t = [v]$ then we have:
    \begin{align*}
      & \llbracket \textsf{let } x_1 \seq x_n \asn ([v_1] \mid \cdots \mid [v_n]) \textsf{ in } [v] \rrbracket
        = ( \llbracket [v_1] \rrbracket \mid \cdots \mid \llbracket [v_n] \rrbracket ) \cdot \llbracket [(v)] \rrbracket
        = \left(\corner{(v_1)} \mid \cdots \mid \corner{(v_n)}\right) \cdot \corner{(v)}
      \\&= \left(\corner{(v_1)} \otimes \cdots \otimes \corner{(v_n)}\right) \cdot \corner{(v)}
      = \corner{(v_1\seq v_n)} \cdot \corner{(v)}
      = \corner{(v_1 \seq v_n)(v)}
      = \corner{(v \circ (v_1 \seq v_n))}
      \\&= \llbracket [v[v_1 \seq v_n/x_1 \seq x_n]] \rrbracket      
    \end{align*}
  \item If $t = \textsf{let } y_1 \seq y_m \asn (t_1 \mid \cdots \mid t_m) \textsf{ in } r$ then we have:
    \begin{align*}
      & \llbracket \textsf{let } x_1^1 \seq x_1^{k_1} \seq x_m^1  \seq x_m^{k_m} \asn ([v_1^1] \mid \cdots \mid [v_1^{k_1}] \mid \cdots \mid [v_m^1] \mid \cdots  \mid [v_m^{k_m}]) \textsf{ in } (\textsf{let } y_1 \seq y_m \asn (t_1 \mid \cdots \mid t_m) \textsf{ in } r) \rrbracket
      \\&= (\llbracket [v_1^1] \rrbracket \mid \cdots \mid \llbracket [v_1^{k_1}]\rrbracket \mid \cdots \mid \llbracket [v_m^{k_m}] \rrbracket \mid \cdots \mid \llbracket [v_m^{k_m}] \rrbracket) \cdot \llbracket \textsf{let } y_1 \seq y_m \asn (t_1 \mid \cdots \mid t_m) \textsf{ in } r) \rrbracket
      \\&= (\llbracket [v_1^1] \rrbracket \mid \cdots \mid \llbracket [v_1^{k_1}]\rrbracket \mid \cdots \mid \llbracket [v_m^{k_m}] \rrbracket \mid \cdots \mid \llbracket [v_m^{k_m}] \rrbracket) \cdot (\llbracket t_1 \rrbracket \mid \cdots \mid \llbracket t_m \rrbracket) \cdot \llbracket r \rrbracket
      \\&= (((\llbracket [v_1^1] \rrbracket \mid \cdots \mid \llbracket [v_1^{k_1}] \rrbracket) \cdot \llbracket t_1 \rrbracket) \mid \cdots \mid ((\llbracket [v_m^1] \rrbracket \mid \cdots \mid \llbracket [v_m^{k_m}] \rrbracket) \cdot \llbracket t_m \rrbracket)) \cdot \llbracket r \rrbracket
      \\&= (\llbracket \textsf{let } x_1^1 \seq x_1^{k_1} \asn ([v_1^1] \mid \cdots \mid [v_1^{k_1}]) \textsf{ in } t_1 \rrbracket \mid \cdots \mid \llbracket \textsf{let } x_m^1 \seq x_m^{k_m} \asn ([v_m^1] \mid \cdots \mid [v_m^{k_m}]) \textsf{ in } t_m \rrbracket) \cdot \llbracket r \rrbracket
      \\&= (\llbracket t_1[v_1^1 \seq v_1^{k_1}/x_1^1 \seq x_1^{k_1}] \rrbracket \mid \cdots \mid \llbracket t_m[v_m^1 \seq v_m^{k_m} / x_m^1 \seq x_m^{k_m}] \rrbracket) \cdot \llbracket r \rrbracket
      \\&= \llbracket \textsf{let } y_1\seq y_m \asn (t_1[v_1^1 \seq v_1^{k_1}/x_1^1\seq x_1^{k_1}] \mid \cdots \mid t_m[v_m^1 \seq v_m^{k_m} / x_m^1 \seq x_m^{k_m}]) \textsf{ in } r \rrbracket
    \end{align*}
  \item If $t = \putr(v,t')$ then we have: 
    \begin{align*}
      & \llbracket \textsf{let } x_1 \seq x_n,y_1\seq y_m \asn ([v_1] \mid \cdots \mid [v_n] \mid [u_1] \mid \cdots \mid [u_m]) \textsf{ in } \putr(v,t')\rrbracket
      \\&= (\llbracket [v_1] \rrbracket \mid \cdots \mid \llbracket [v_n] \rrbracket \mid \llbracket [u_1] \rrbracket \mid \cdots \mid \llbracket [u_m] \rrbracket) \cdot \llbracket \putr(v,t') \rrbracket
      \\&= (\llbracket [v_1] \rrbracket \mid \cdots \mid \llbracket [v_n] \rrbracket \mid \llbracket [u_1] \rrbracket \mid \cdots \mid \llbracket [u_m] \rrbracket) \cdot (1 \mid (\corner{(v)} \cdot A\llcorner)) \cdot \llbracket t' \rrbracket
      \\&= ((\llbracket [v_1] \rrbracket \mid \cdots \llbracket [v_n] \rrbracket) \cdot \llbracket t' \rrbracket) \mid ((\llbracket [u_1] \rrbracket \mid \cdots \mid \llbracket [u_m] \rrbracket) \cdot \corner{(v)} \cdot A\llcorner)
      \\&= \llbracket t'[v_1 \seq v_n/x_1 \seq x_n] \rrbracket \mid (\corner{( v[u_1\seq u_m / y_1 \seq y_m] )} \cdot A\llcorner)
      \\&= (1 \mid (\corner{( v[u_1\seq u_m / y_1 \seq y_m] )} \cdot A\llcorner)) \mid \llbracket t'[v_1 \seq v_n / x_1 \seq x_n] \rrbracket 
      \\&= \llbracket \putr(v[u_1\seq u_m/y_1\seq y_m],t'[v_1 \seq v_n/x_1 \seq x_n]) \rrbracket
      \\&= \llbracket \putr(v,t')[v_1\seq v_n,u_1 \seq u_m/x_1 \seq x_n,y_1 \seq y_m] \rrbracket
    \end{align*}
  \item If $t = \putl(v,t')$, a similar argument to the one above suffices.
  \item If $t = \getr(x.t')$, then we have:
    \begin{align*}
      & \llbracket \textsf{let } x_1 \seq x_n \asn ([v_1] \mid \cdots \mid [v_n]) \textsf{ in } \getr(x.t') \rrbracket
        = (\llbracket [v_1] \rrbracket \mid \cdots \mid \llbracket [v_n] \rrbracket ) \cdot \llbracket \getr(x.t') \rrbracket
      \\&= (\llbracket [v_1] \rrbracket \mid \cdots \mid \llbracket [v_n] \rrbracket ) \cdot (1 \mid A\ulcorner) \cdot \llbracket t' \rrbracket
      = \corner{(v_1\seq v_n)} \cdot (1 \mid A\ulcorner) \cdot \llbracket t' \rrbracket
      \\&= (1 \mid A\ulcorner) \cdot \corner{(v_1 \seq v_n,1_A)} \cdot \llbracket t' \rrbracket
      = (1 \mid A\ulcorner) \cdot \llbracket \textsf{let } x_1 \seq x_n,x \asn ([v_1] \mid \cdots [v_n] \mid [x]) \textsf{ in } t' \rrbracket
      \\&= (1 \mid A\ulcorner) \cdot \llbracket t'[v_1\seq v_n,x / x_1 \seq x_n,x] \rrbracket
      = \llbracket \getr(x.t'[v_1 \seq v_n,x / x_1 \seq x_n,x]) \rrbracket
      \\&= \llbracket \getr(x.t')[v_1 \seq v_n / x_1 \seq x_n]\rrbracket
    \end{align*}
  \item If $t = \getl(x.t')$, a similar argument to the one above suffices.
  \end{itemize}
\end{proof}

\end{document}
